\input{../tex-inputs/latex-leseansicht-vorspann}

\section[Arthur Schnitzler an Gustav Schwarzkopf, 3. 4. 1901]{L04414 Arthur Schnitzler an Gustav Schwarzkopf, 3. 4. 1901}
\nopagebreak\mylabel{L04414v}
\rehead{ }\normalsize\beginnumbering\briefempfaengerindex{Schwarzkopf, Gustav@\textsc{Schwarzkopf, Gustav}!zzzSchnitzler, Arthur@\emph{von Arthur Schnitzler}!1901-04-032@{3. 4. 1901}|(be}
\toendnotes[C]{\smallbreak\pagebreak[2]}
\correspDesc{Versand  durch Arthur Schnitzler am 3. 4. 1901 in Rom
\newline{}Übermittlung  am 3. 4. 1901 in Rom
\newline{}Zustellung  am 5. 4. 1901 in Wien
\newline{}Erhalt  durch Gustav Schwarzkopf am 5. 4. 1901 in Wien}\toendnotes[C]{\smallbreak}
\Standort{DLA, A:Schnitzler, HS.1985.1.1897.}
\physDesc{Bildpostkarte, 120 Zeichen
\newline{}Handschrift: Bleistift, deutsche Kurrent
\newline{}Versand: 1) Stempel: »\nobreak{}\oindex{Rom@\textbf{Rom}|pwk}Ro\textcolor{gray}{m}a \textcolor{gray}{Ferrovi}a, 3 \textcolor{gray}{4}–0\textcolor{gray}{1}\nobreak{}«.   2) Stempel: »\nobreak{}\oindex{I., Innere Stadt@\textbf{I., Innere Stadt}|pwk}Wien 1/1 1, 5. 4. \textcolor{gray}{1}901, 8–9½\textcolor{gray}{V}., Bestellt\nobreak{}«. }\toendnotes[C]{\smallbreak}\pstart{}{\pb}Hrn Guſtav Schwarzkopf\pend{}\pstart{}Wien\oindex{Wien@\textbf{Wien}|pw}\pend{}\pstart{}\textsc{I. Tiefer Graben 23}\oindex{Wien@\textbf{Wien}!I., Innere Stadt@\textbf{I., Innere Stadt}!Tiefer Graben 23@\textbf{Tiefer Graben 23}|pw}.\pend{}{\bigskip}
\pstart
           \centering{}{\pb}\textcolor{gray}{\textbf{41. Foro Romano – Tempio di
                     Castore e Polluce\oindex{Tempio dei Càstori@\textbf{Tempio dei Càstori}|pw}.}}\pend
           \vspace{1em}
\pstart
           \noindent{}{\pb}Herzliche Grüße! Warum \label{K_L04414-1v}\edtext{intriguiren}{\lemma{\textnormal{\emph{intriguiren}}}\Cendnote{\textnormal{unklare Anspielung; wahrscheinlich scherzhaft.}}}\label{K_L04414-1} Sie de{\geminationn} gegen
               mich!? –\pend
           \pstart Ihr \spacefill\mbox{Arthur}\pend{}
\pstart
           3. 4. 901\pend
           \selectlanguage{ngerman}\endnumbering\briefempfaengerindex{Schwarzkopf, Gustav@\textsc{Schwarzkopf, Gustav}!zzzSchnitzler, Arthur@\emph{von Arthur Schnitzler}!1901-04-032@{3. 4. 1901}|)be}\mylabel{L04414h}
\begin{anhang}
\end{anhang}\newcommand{\dateiname}{L04414}\newcommand{\titel}{Arthur Schnitzler an Gustav Schwarzkopf, 3. 4. 1901}\newcommand{\editorInnen}{Herausgegeben von Jahnke, SelmaMüller, Martin Anton}\input{../tex-inputs/latex-leseansicht-abspann}
